\documentclass[10pt,a4paper]{article}
\usepackage[a4paper,hmargin=16mm,top=13mm,bottom=16mm]{geometry}
\usepackage{lmodern}
\usepackage[T1]{fontenc}
\renewcommand{\familydefault}{\sfdefault}
\usepackage{pifont,longtable,booktabs,array,needspace}
\usepackage[hidelinks]{hyperref}
\setlength{\parindent}{0pt}
\setlength{\parskip}{4pt}
\setlength{\LTpre}{2pt}
\setlength{\LTpost}{4pt}
\renewcommand{\arraystretch}{1.12}

% ---- Checkboxes: \cb is open, \cbdone is ticked. Each is a clickable PDF field.
\newcounter{cbn}
\newcommand{\cb}{\mbox{\stepcounter{cbn}\CheckBox[name=cb\thecbn,width=9pt,height=9pt,borderwidth=0.8pt,bordercolor=0 0 0,backgroundcolor=1 1 1,checkboxsymbol=\ding{51}]{}}}
\newcommand{\cbdone}{\mbox{\stepcounter{cbn}\CheckBox[name=cb\thecbn,checked=true,width=9pt,height=9pt,borderwidth=0.8pt,bordercolor=0 0 0,backgroundcolor=1 1 1,checkboxsymbol=\ding{51}]{}}}

% ---- Daily plan rows
\newenvironment{days}{\begin{longtable}{@{}p{6mm}p{9mm}p{154mm}@{}}}{\end{longtable}}
\newcommand{\dayrow}[1]{\midrule\multicolumn{3}{@{}l@{}}{\textbf{#1}}\\*}
\newcommand{\task}[2]{\cb & #1 & #2 \\}
\newcommand{\taskdone}[2]{\cbdone & #1 & #2 \\}
\newcommand{\review}{\textbf{Sunday review:}\quad \cb~on track\qquad \cb~behind\qquad Recovery plan and owner: \hrulefill\par\medskip}
\newcommand{\week}[2]{\needspace{9\baselineskip}\section*{#1}\vspace{-6pt}\textit{Goal: #2}\par}

\begin{document}
\begin{Form}

\begin{center}
{\LARGE\bfseries NeuroPresence -- FYP-I Progress Plan}\\[4pt]
Daily checklist, 6 October to 7 December 2026\\[2pt]
Zain Shahid (Z) \quad Muhammad Talha Arshad (T) \quad Sana Ullah Farooqi (S) \quad Supervisor: Muhammad Aamir Gulzar
\end{center}
\hrule\medskip

\textbf{How to use this sheet.}
Tick a box only when the work is pushed to GitHub and can be shown running or opened; a verbal ``done'' does not count.
Boxes are clickable in Adobe Reader or Microsoft Edge (save the file after ticking), or print the sheet and tick by hand.
Every Sunday, mark the week on track or behind and write the recovery plan with an owner.
``All'' means the three members together. Items marked \emph{(stretch)} are dropped first if the week runs late.

\textbf{Dates this plan assumes.}

\begin{tabular}{@{}p{50mm}p{34mm}p{80mm}@{}}
\toprule
Event & Date & Basis \\
\midrule
FYP-I mid evaluation & Wed 28 Oct 2026 & Assumed by the team; change if the supervisor sets another date \\
FYP-I final evaluation & Mon 7 Dec 2026 & Not announced; placeholder until the department publishes it \\
\bottomrule
\end{tabular}

\textbf{Open decisions to close this week.}

\begin{tabular}{@{}p{6mm}p{166mm}@{}}
\cbdone & Development platform: both. Windows on Zain's PC and Kubuntu on a teammate's machine; all code must run on both. \\
\cb & How Talha and Sana get working time on the RTX 5050 machine (schedule or remote access). \\
\cb & Repository visibility: it is public now; decide whether it stays public. \\
\end{tabular}

\section*{Milestones}

\begin{longtable}{@{}p{6mm}p{10mm}p{118mm}p{28mm}@{}}
\toprule
 & No. & Milestone & Due \\
\midrule
\endhead
\cbdone & M1 & Repository created; capture and tracking module with tests & Tue 6 Oct \\
\cb & M2 & Baseline reproduced from the repository (7.7 fps, 130 ms, 1.57 GB) & Sun 11 Oct \\
\cb & M3 & Live reenactment: webcam to tracker to LivePortrait to virtual camera & Sun 18 Oct \\
\cb & M4 & Mid report, diagrams, test evidence, and action register complete & Sun 25 Oct \\
\cb & M5 & \textbf{Mid evaluation} & Wed 28 Oct \\
\cb & M6 & Robust core pipeline; optimization findings recorded & Sun 8 Nov \\
\cb & M7 & SRS and SDS baselines with the supervisor (Gates 2 and 3) & Sun 15 Nov \\
\cb & M8 & Clean-machine install and executed test plan; feature freeze & Sun 22 Nov \\
\cb & M9 & FYP-I final report draft with the supervisor & Sun 29 Nov \\
\cb & M10 & \textbf{FYP-I final evaluation}: working core system (Final-1) & Mon 7 Dec \\
\bottomrule
\end{longtable}

\section*{Mid evaluation deliverables: demo and presentation (Form 2, 100 marks)}

\begin{longtable}{@{}p{6mm}p{44mm}p{11mm}p{90mm}p{9mm}@{}}
\toprule
 & Criterion & Marks & What we show & Who \\
\midrule
\endhead
\cb & 1. Proposal actions addressed & 10 & Panel Action Register: every change the proposal panel asked for, its status, and a link to the evidence & S \\
\cb & 2. Iteration plan and progress & 10 & This sheet and the task log, with an owner per task and a recovery plan for anything late & All \\
\cb & 3. Live core module demo & 25 & Webcam to tracking to reenactment to virtual camera, running live; then no face, two faces, face turned away, camera unplugged & All \\
\cb & 4. Requirements quality & 10 & Numbered FRs and NFRs, each with a testable acceptance check and a pointer to the code or test & S \\
\cb & 5. Design artefacts and evolution & 15 & Architecture, sequence, data-flow, and state diagrams that match the code, with changes since the proposal explained & All \\
\cb & 6. Repository and member evidence & 10 & Regular commits from all three members with meaningful messages; tagged mid release & All \\
\cb & 7. Testing and early validation & 10 & Tests run live with \texttt{pytest}; edge-case log; one failure case discussed & All \\
\cb & 8. Presentation and ownership & 10 & Each member explains their own module at function level and defends any tool-assisted code & All \\
\bottomrule
\end{longtable}

\section*{Mid evaluation deliverables: mid report (Form 3, 100 marks)}

\begin{longtable}{@{}p{6mm}p{52mm}p{11mm}p{82mm}p{9mm}@{}}
\toprule
 & Section & Marks & Content & Who \\
\midrule
\endhead
\cb & 1. Format and artefacts & 5 & Prescribed template, all required sections, signatures & S \\
\cb & 2. Problem, R\&D, and gap analysis & 15 & Literature synthesis written in our own words; the gap and the complex computing challenge isolated & Z \\
\cb & 3. Vision, scope, and success criteria & 10 & Scope table, module breakdown, measurable success criteria, consistent with the proposal & S \\
\cb & 4. Requirements and acceptance criteria & 15 & Numbered FRs and NFRs with verifiable acceptance checks & S \\
\cb & 5. Architecture and interface design & 15 & Layered architecture and component interfaces with data shapes, justified & T \\
\cb & 6. Data, process, and behaviour design & 10 & Data-flow, sequence, and state diagrams, each explained in the text & Z \\
\cb & 7. Implementation documentation & 15 & Reproducible setup, main files, classes, functions, and technical decisions & T \\
\cb & 8. Test plan, tools, and risks & 15 & Test plan and early results; every external model, licence, and GenAI tool disclosed; risks & All \\
\bottomrule
\end{longtable}

\newpage
% =====================================================================
\week{Week 1: Tue 6 Oct -- Sun 11 Oct. Setup and baseline}{everyone can run the code, and the proposal baseline is reproduced from the repository.}

\begin{days}
\dayrow{Tue 6 Oct}
\taskdone{Z}{Create the GitHub repository with README and \texttt{.gitignore}}
\taskdone{Z}{Capture and face-tracking module (stages 1--2) with 10 passing tests}
\task{Z}{Push the two local commits to GitHub}
\task{All}{Add Talha, Sana, and the supervisor as collaborators}
\dayrow{Wed 7 Oct}
\task{Z}{Run the capture demo on a real webcam; confirm landmarks and pose angles look right}
\taskdone{T}{Clone the repository; install GPU PyTorch, LivePortrait, and its weights on the RTX 5050 machine}
\task{S}{Collect the proposal-defense feedback into a Panel Action Register in \texttt{docs/}}
\dayrow{Thu 8 Oct}
\task{T}{Commit the proposal benchmark script and \texttt{results/real\_run.json}; reproduce the baseline}
\task{Z}{Record the source clip and two driving test clips of each member (kept out of the repository)}
\task{S}{Get the mid-report template and FYP handbook Forms 2, 3, 10, and 11; create the report skeleton}
\dayrow{Fri 9 Oct}
\taskdone{T}{LivePortrait offline: one source image plus a driving video file produces an output video}
\taskdone{Z}{Crop and align the tracked face to the reenactment input size}
\task{S}{Task log in the repository with one entry per task, owner, and status}
\dayrow{Sat 10 Oct}
\taskdone{T}{Split LivePortrait into ``prepare source once'' and ``drive one frame'' functions}
\taskdone{Z}{Tests for the crop, including a face at the frame edge}
\task{S}{SRS draft: number the FYP-I requirements and write one acceptance check for each}
\dayrow{Sun 11 Oct}
\task{All}{Confirm every member has pushed at least one commit this week}
\task{All}{Run the test suite on both Windows and Kubuntu; fix anything that passes on only one}
\end{days}
\review

% =====================================================================
\week{Week 2: Mon 12 Oct -- Sun 18 Oct. Live reenactment loop}{the full path from webcam to a meeting application runs live.}

\begin{days}
\dayrow{Mon 12 Oct}
\taskdone{T}{Live loop: webcam frame to LivePortrait to a preview window (first live reenactment)}
\taskdone{Z}{Gate the loop on tracker status: reenact only when exactly one face is found}
\task{S}{Metrics module (FR-7): per-stage latency, fps, and peak VRAM, on screen and logged to file}
\dayrow{Tue 13 Oct}
\task{T}{Measure per-stage latency in the live loop and commit the numbers}
\task{Z}{Smoothing filter on the driving signal, with an on/off switch}
\task{S}{Virtual camera writer, one code path: OBS Virtual Camera on Windows, v4l2loopback on Kubuntu}
\dayrow{Wed 14 Oct}
\taskdone{T}{Fallback output when there is no usable face: hold the static enrolled frame}
\task{Z}{Frame pacing: fixed deadline, drop late frames, never block}
\task{S}{Test the virtual camera in Google Meet; start the compatibility matrix}
\dayrow{Thu 15 Oct}
\task{T}{Tests for the reenactment wrapper: output shape, source caching, bad input}
\task{Z}{Tests for pacing and smoothing}
\task{S}{Test the virtual camera in the Zoom desktop client; record client versions}
\dayrow{Fri 16 Oct}
\task{All}{Integration: webcam to tracker to reenactment to virtual camera to Google Meet, end to end}
\task{S}{Update the architecture diagram to match the code as built}
\dayrow{Sat 17 Oct}
\task{T}{Ten-minute stability run; log latency and VRAM; fix crashes}
\task{Z}{Edge-case log: no face, two faces, face turned away, low light, camera unplugged}
\task{S}{Install ArcFace; measure CSIM between source and output on a recorded run}
\dayrow{Sun 18 Oct}
\task{All}{Record a short capture of the working pipeline as a reference (not a substitute for the live demo)}
\end{days}
\review

% =====================================================================
\week{Week 3: Mon 19 Oct -- Sun 25 Oct. Design evidence and mid report}{the report and diagrams are complete and match the code.}

\begin{days}
\dayrow{Mon 19 Oct}
\task{S}{Sequence diagram of one frame through the pipeline; data-flow diagram}
\task{Z}{State diagram: tracking OK, no face, multiple faces, fallback}
\task{T}{Interface description for the reenactment component: inputs, outputs, tensor shapes}
\dayrow{Tue 20 Oct}
\task{Z}{Report section 2: problem, R\&D, and gap analysis, in our own words}
\task{S}{Report section 3: vision, scope table, and success criteria}
\task{T}{Report section 7: setup steps, main files, classes, and functions}
\dayrow{Wed 21 Oct}
\task{Z}{Fix the proposal inconsistencies in the report: watermark optional or always-on; scope modes and languages; fallback type; bandwidth claim}
\task{S}{Report section 4: requirements traced to code and tests}
\task{T}{Follow the setup steps from a clean clone and correct anything that fails}
\dayrow{Thu 22 Oct}
\task{T}{Report section 5: architecture and component interfaces}
\task{Z}{Report section 6: explain each diagram in the text}
\task{S}{Report section 8: test plan and early results table (latency, fps, VRAM, CSIM)}
\dayrow{Fri 23 Oct}
\task{All}{Tool and dependency disclosure: every external model, licence, and GenAI tool, and where each was used}
\task{S}{Panel Action Register: status and evidence link for every action}
\task{All}{Update the risk table with what the first three weeks showed}
\dayrow{Sat 24 Oct}
\task{All}{Assemble the report in the template; each member reads the other two members' sections}
\task{All}{Send the report draft to the supervisor}
\dayrow{Sun 25 Oct}
\task{All}{Draft the presentation slides and the demo script}
\end{days}
\review

% =====================================================================
\week{Week 4: Mon 26 Oct -- Sun 1 Nov. Rehearsal, mid evaluation, feedback}{a clean live demo, then every piece of feedback turned into a task.}

\begin{days}
\dayrow{Mon 26 Oct}
\task{All}{Rehearsal 1: happy path and four failure cases, timed}
\task{All}{Each member explains their own module line by line to the other two}
\task{All}{Fix what the rehearsal exposed}
\dayrow{Tue 27 Oct}
\task{All}{Rehearsal 2 on the actual demo machine with the actual webcam}
\task{S}{Final report PDF and slides ready; task log up to date}
\task{T}{Tag the release \texttt{v0.1-mid}; confirm the commit history shows all three members}
\dayrow{Wed 28 Oct -- MID EVALUATION}
\task{All}{Demo, presentation, and report evaluated}
\task{S}{Write down every supervisor comment the same day}
\dayrow{Thu 29 Oct}
\task{All}{Enter each comment in the Panel Action Register with an owner and a due date}
\task{All}{Adjust the rest of this plan for anything the feedback changes}
\dayrow{Fri 30 Oct}
\task{T}{Profile the live loop per module and rank the three largest costs}
\task{Z}{Script to measure inter-frame warping error on a recorded run (temporal-stability baseline)}
\task{S}{Move the source clip path, thresholds, and camera index into a config file}
\dayrow{Sat 31 Oct}
\task{T}{Persistent FP16 weights for the warping module and generator; re-measure}
\task{Z}{Compare smoothing on and off with the warping-error script}
\task{S}{Benchmark harness v1: fixed driving clips in; latency, fps, VRAM, and CSIM out as JSON}
\dayrow{Sun 1 Nov}
\task{All}{Confirm the feedback tasks all have owners and dates}
\end{days}
\review

% =====================================================================
\week{Week 5: Mon 2 Nov -- Sun 8 Nov. Robustness and optimization findings}{the pipeline survives real conditions, and we know what moves latency.}

\begin{days}
\dayrow{Mon 2 Nov}
\task{T}{\texttt{torch.compile} on the warping module; record latency before and after}
\task{Z}{Camera-loss recovery: detect an unplugged camera, retry, and show a clear message}
\task{S}{Run the harness on the unoptimized baseline with team clips; commit the results}
\dayrow{Tue 3 Nov}
\task{T}{\texttt{torch.compile} on the generator; record latency before and after}
\task{Z}{Find the yaw and pitch at which the output breaks; hold the last good frame beyond them}
\task{S}{Complete the compatibility matrix for Meet, Zoom, and Teams on Windows and Kubuntu, with client versions}
\dayrow{Wed 4 Nov}
\task{T}{Resolution scheduling experiment: latency against image quality}
\task{Z}{Tests for camera-loss and pose-limit handling}
\task{S}{Pass-through fallback: switch to the real camera if reenactment fails or exceeds the budget}
\dayrow{Thu 5 Nov}
\task{T}{Findings note in \texttt{docs/}: what reduced latency and what did not}
\task{Z}{Measure the offset between voice and reenacted lips in a recorded call}
\task{S}{Tests for the fallback switch and config loading}
\dayrow{Fri 6 Nov}
\task{All}{Merge all branches; run the full pipeline for 15 minutes; log every failure}
\dayrow{Sat 7 Nov}
\task{All}{Fix the failures from the 15-minute run}
\task{S}{Update the task log and milestone status}
\dayrow{Sun 8 Nov}
\task{All}{Decide which stretch items for Week 6 are kept}
\end{days}
\review

% =====================================================================
\week{Week 6: Mon 9 Nov -- Sun 15 Nov. SRS and SDS baselines}{requirements and design documents are with the supervisor (Gates 2 and 3).}

\begin{days}
\dayrow{Mon 9 Nov}
\task{S}{SRS: all FRs and NFRs final, each marked FYP-I or FYP-II, each with an acceptance check}
\task{Z}{SRS: user interactions (enrol the source clip, start a session, fallback prompts)}
\task{T}{SRS: how each target is measured (42 ms, 24 fps, 8 GB, CSIM 0.80)}
\dayrow{Tue 10 Nov}
\task{T}{SDS: reenactment component design, interfaces, and data shapes}
\task{Z}{SDS: capture, tracking, and pacing design; state diagram updated}
\task{S}{SDS: architecture, deployment view, virtual camera, and safeguard design}
\dayrow{Wed 11 Nov}
\task{S}{\emph{(stretch)} Enrolment: save the enrolled face embedding; self-likeness check with a threshold}
\task{T}{\emph{(stretch)} ONNX export of the motion extractor; verify its output matches PyTorch}
\task{Z}{Update the sequence diagram to match the merged code}
\dayrow{Thu 12 Nov}
\task{S}{Always-on corner watermark on the output stream (FR-5)}
\task{T}{Live CSIM readout every few frames in the metrics overlay}
\task{Z}{Check every diagram against the code and fix mismatches}
\dayrow{Fri 13 Nov}
\task{All}{Send the SRS and SDS drafts to the supervisor}
\dayrow{Sat 14 Nov}
\task{T}{Thirty-minute endurance trial: log VRAM, latency, and CSIM over time (formal run stays in FYP-II)}
\task{Z}{Tests for the watermark}
\task{S}{Tests for the enrolment check, if built}
\dayrow{Sun 15 Nov}
\task{All}{Confirm the FYP-I final date and rubric with the supervisor; update this sheet}
\end{days}
\review

% =====================================================================
\week{Week 7: Mon 16 Nov -- Sun 22 Nov. Clean install, test execution, freeze}{someone who did not write the code can install and run it; features stop changing.}

\begin{days}
\dayrow{Mon 16 Nov}
\task{T}{Documented install steps for Windows and Kubuntu with pinned versions}
\task{Z}{Test plan document: unit, integration, and edge cases with expected results}
\task{S}{Install on a clean environment by a member who did not write the steps; record every gap}
\dayrow{Tue 17 Nov}
\task{T}{Fix the install gaps}
\task{Z}{Execute the test plan; record a pass/fail table}
\task{S}{Run the harness on the current pipeline; results table for the report}
\dayrow{Wed 18 Nov}
\task{T}{Fix defects from the test run (reenactment)}
\task{Z}{Fix defects from the test run (capture and pacing)}
\task{S}{Fix defects from the test run (output, fallback, config)}
\dayrow{Thu 19 Nov}
\task{All}{Apply the supervisor's feedback to the SRS and SDS}
\dayrow{Fri 20 Nov}
\task{All}{Feature freeze: only bug fixes from here; tag \texttt{v0.2-rc}}
\dayrow{Sat 21 Nov}
\task{Z}{Final report: update problem, R\&D, and gap analysis}
\task{T}{Final report: implementation and measured results}
\task{S}{Final report: requirements, design, and testing}
\dayrow{Sun 22 Nov}
\task{All}{Re-run the full test suite on the frozen code}
\end{days}
\review

% =====================================================================
\week{Week 8: Mon 23 Nov -- Sun 29 Nov. Final report and documents}{a complete FYP-I report draft is with the supervisor.}

\begin{days}
\dayrow{Mon 23 Nov}
\task{Z}{Final report: data, process, and behaviour design with final diagrams}
\task{T}{Final report: architecture and optimization findings}
\task{S}{Final report: scope, success criteria, and what moved to FYP-II}
\dayrow{Tue 24 Nov}
\task{All}{Final report: test results, risks, tool and licence disclosure}
\task{S}{Panel Action Register: close every item with evidence}
\dayrow{Wed 25 Nov}
\task{All}{Each member reads the full report and marks anything they cannot explain}
\dayrow{Thu 26 Nov}
\task{All}{Rewrite the marked passages; check every number against the committed results}
\dayrow{Fri 27 Nov}
\task{All}{Send the full report draft to the supervisor}
\dayrow{Sat 28 Nov}
\task{All}{Draft the final presentation slides and demo script}
\dayrow{Sun 29 Nov}
\task{All}{Write the FYP-II plan for the first four weeks of Spring 2027}
\end{days}
\review

% =====================================================================
\week{Week 9: Mon 30 Nov -- Mon 7 Dec. Rehearsal and FYP-I final}{a rehearsed live demo and submitted documents.}

\begin{days}
\dayrow{Mon 30 Nov}
\task{All}{Apply the supervisor's feedback to the report}
\dayrow{Tue 1 Dec}
\task{All}{Rehearsal 1: full demo with failure cases, timed}
\dayrow{Wed 2 Dec}
\task{All}{Fix what the rehearsal exposed (bug fixes only)}
\dayrow{Thu 3 Dec}
\task{All}{Rehearsal 2 with mock panel questions; each member answers on their own module}
\dayrow{Fri 4 Dec}
\task{S}{Submit the final report, SRS, and SDS}
\task{T}{Tag the release \texttt{v1.0-fyp1}}
\dayrow{Sat 5 Dec}
\task{All}{Dry run on the demo machine; check webcam, virtual camera, and meeting client}
\dayrow{Sun 6 Dec}
\task{All}{No changes to code; pack the demo equipment}
\dayrow{Mon 7 Dec -- FYP-I FINAL EVALUATION (assumed date)}
\task{All}{Panel defense and live demo}
\task{S}{Record the panel's actions for FYP-II}
\end{days}

\section*{FYP-I final deliverables}

The final rubric has not been released; this list follows the proposal timeline and should be updated when it is.

\begin{longtable}{@{}p{6mm}p{124mm}p{38mm}@{}}
\toprule
 & Deliverable & Who \\
\midrule
\endhead
\cb & Working core system: live webcam to reenactment to virtual camera in a meeting application (Final-1) & All \\
\cb & Live metrics shown to the user: latency, fps, and GPU memory (FR-7) & S \\
\cb & Fallbacks working: static enrolled frame, and pass-through camera on failure & T, S \\
\cb & Always-on disclosure watermark on the output (FR-5) & S \\
\cb & Measured results committed: baseline, optimization experiments, CSIM, temporal stability & T, Z, S \\
\cb & SRS baseline (Gate 2) and SDS baseline (Gate 3) & All \\
\cb & FYP-I final report in the prescribed template & All \\
\cb & Test plan with executed results & Z \\
\cb & Meeting-client compatibility matrix for Windows and Kubuntu & S \\
\cb & Install steps verified on a clean environment on both Windows and Kubuntu & T \\
\cb & Panel Action Register closed with evidence & S \\
\cb & Tagged release \texttt{v1.0-fyp1} with commits from all three members & All \\
\bottomrule
\end{longtable}

\textbf{Deferred to FYP-II, as in the proposal:} reaching the 42 ms budget with TensorRT, the full consent gate and identity fallback, the offline audio and text mode, the live audio-driven mode with the phone relay, Urdu voice adaptation, the full benchmark, and the user-acceptance study.

\end{Form}
\end{document}
